\renewcommand{\Anno}{2012}

\renewcommand{\NumeroProblema}{1}

\renewcommand{\Testo}{
L'interno della vasca di un acquario ha la forma di un
parallelepipedo.  Sul fondo dell'acquario \`{e} appoggiato un prisma
(si intende che l'acquario \`{e} a sua volta appoggiato su un piano
orizzontale e la base del prisma \`{e} totalmente a contatto con la
base dell'interno dell'acquario).  Non sono note le misure interne
della vasca, n\'{e} l'altezza del prisma, n\'{e} la forma o le misure
della base del prisma.

Inizialmente l'acquario \`{e} vuoto.  Successivamente, Paolo versa 10
litri di acqua nell'acquario, e nota che il livello sale a 12~cm.  Poi
versa altri 10 litri di acqua, ed il livello sale a 22~cm.  Infine
versa altri 10 litri di acqua, ed il livello sale a 30~cm (tutti i 
livelli sono riferiti alla parte interna del fondo della vasca).

\begin{enumerate}
	\renewcommand{\labelenumi}{(\alph{enumi})}
	\item  Determinare quale livello raggiungerebbe l'acqua se 
	Paolo versasse altri 10 litri di acqua (supponendo ovviamente che 
	ci stiano nell'acquario). 

	\item Determinare quanti litri di acqua avrebbe dovuto versare
	inizialmente Paolo per fare arrivare l'acqua esattamente in
	corrispondenza della base superiore del prisma.
\end{enumerate}
}

\renewcommand{\Risposta}{
\begin{enumerate}
	\renewcommand{\labelenumi}{(\alph{enumi})}
	\item  38 cm.

	\item  15 litri
\end{enumerate}
}

\renewcommand{\Aiutino}{
Durante quale dei 3 versamenti di Paolo avviene che l'acqua ricopre
completamente il prisma?  Di quanto cresce il livello dell'acqua dopo
il superamento della sommit\`{a} del prisma?  Di quanto cresce il
livello dell'acqua prima del superamento della sommit\`{a} del 
prisma? Quale o quali formule descrivono il volume dell'acqua in 
funzione del livello raggiunto?}

\renewcommand{\SoluzioneUfficiale}{\TitoloSoluzione 
Sia $h$ l'altezza del prisma. Fino a quando non raggiunge la 
sommit\`{a} del prisma, il livello dell'acqua cresce linearmente con 
il volume versato, con una costante di proporzionalit\`{a} dipendente 
dalla differenza tra l'area base $S_{V}$ della vasca e l'area base 
$S_{P}$ del prisma. In formule, il volume $V$ ed il livello $\ell$ 
sono legati dalla relazione
\begin{equation}
	V=(S_{V}-S_{P})\ell
	\hspace{3em}
	0\leq\ell\leq h.
	\label{eqn:vol-sotto}
\end{equation} 

Non appena il livello supera la quota $h$, la costante di 
proporzionalit\`{a} dipende solo pi\`{u} da $S_{V}$, e si ha che
\begin{equation}
	V=(S_{V}-S_{P})h+S_{V}(\ell-h)
	\hspace{3em}
	\ell\geq h.
	\label{eqn:vol-sopra}
\end{equation}

Dai dati del problema \`{e} chiaro quindi che la sommit\`{a} del 
prisma \`{e} stata raggiunta tra il decimo ed il ventesimo litro. Da 
quel punto in poi, la crescita del livello \`{e} proporzionale 
all'aumento di volume, e la costante di proporzionalit\`{a} non 
cambia pi\`{u}. Se dunque il livello sale di 8~cm passando da 20 a 30 
litri, esso salir\`{a} di 8~cm anche passando da 30 a 40 litri. 
Pertanto, se si versassero 40 litri di acqua, il livello totale 
sarebbe di 38~cm, il che risponde al punto~(a).

Per rispondere al punto~(b), indichiamo con $x$ il numero di litri
richiesto, che sappiamo gi\`{a} essere compreso tra 10 e 20.
Dai dati del problema, deduciamo che prima del raggiungimento della
sommit\`{a} del prisma il livello cresce di 12~cm ogni 10 litri di
acqua versati, e dopo cresce di 8~cm ogni 10 litri di acqua versati. 
Pertanto, dopo 20 litri di acqua il livello in cm sar\`{a}
$$22=\frac{12}{10}x+\frac{8}{10}(20-x),$$
da cui $x=15$.
}

\renewcommand{\SoluzioniAlternative}{\TitoloSoluzioneAlternativaA
Usando le equazioni (\ref{eqn:vol-sotto}) e (\ref{eqn:vol-sopra}) per 
i 3 valori dati di volume (tradotto in cm$^{3}$) e livello (in cm), 
si ottiene il sistema
$$\left\{
\begin{array}{l}
	(S_{V}-S_{P})12=10.000,  \\
	\noalign{\vspace{0.5ex}}
	(S_{V}-S_{P})h+S_{V}(22-h)=20.000,  \\
	\noalign{\vspace{0.5ex}}
	(S_{V}-S_{P})h+S_{V}(30-h)=30.000,
\end{array}
\right.$$
da cui con semplici calcoli si ricavano le 3 incognite
$$h=18\mbox{ cm},
\quad\quad
S_{V}=1.250\mbox{ cm}^{2},
\quad\quad
S_{P}=\frac{1.250}{3}\mbox{ cm}^{2}.$$

Da qui \`{e} facile rispondere ai punti~(a) e~(b). Per il punto~(a), 
basta applicare la (\ref{eqn:vol-sopra}) con $V=40.000$. Si ottiene 
cos\`{i} che
$$40.000=\left(1.250-\frac{1.250}{3}\right)\cdot 18+
1.250(\ell-18),$$
da cui con semplici calcoli si ottiene $\ell=38$.

Per il punto~(b), basta applicare la (\ref{eqn:vol-sotto}) o la 
(\ref{eqn:vol-sopra}), che in questo caso sono equivalenti, con 
$\ell=18$. Si ottiene cos\`{i} che
$$V=\left(1.250-\frac{1.250}{3}\right)\cdot 18=15.000,$$
che corrisponde a 15 litri.
}

\renewcommand{\Osservazioni}{}

\renewcommand{\Commenti}{
}